% GENERATED FILE. Edit canonical lessons, then run npm run generate:books.
% canonical-source-hash: fnv1a64:5e7b15a2a89ed52a
% canonical-lessons: SA-C146-pautrah, SA-C146-pautri, SA-C146-pitamahi, SA-C146-matamahah, SA-C146-pitrvyah

\chapter{Grandson, Granddaughter, Paternal grandmother, Maternal grandfather, Father's brother}
\label{ch:sa-grandson-granddaughter-paternal-grandmother-maternal-grandfather-father-s-brother}
\hlchaptermodality{146}

\begin{chapteropening}
\textbf{By the end of this chapter:} \emph{I can say a grandson, a granddaughter, a paternal grandmother, a maternal grandfather and a father's brother.}

Complete the last lesson of chapter 146: I can say a grandson, a granddaughter, a paternal grandmother, a maternal grandfather and a father's brother.
\end{chapteropening}

\section[\texorpdfstring{\sk{पौत्रः} (pautraḥ)}{पौत्रः (pautraḥ)}]{\sk{पौत्रः} (pautraḥ) — a grandson}

\label{lesson:SA-C146-pautrah}



\begin{quote}
Before the new one: say the Sanskrit for coos, then the Sanskrit for chatters.
\end{quote}

\subsection*{You'll want to know: \sk{पौत्रः}}
\textbf{\sk{पौत्रः}} — \emph{pautraḥ} — \textquotedblleft{}a grandson\textquotedblright{}.

\begin{grammarlens}[title={What you've built}]
One more word for this chapter.
\end{grammarlens}

\subsection*{Your turn}
\begin{itemize}
\raggedright
  \item \emph{Say it:} \emph{pautraḥ}
  \item \emph{Say it:} \emph{pautraḥ}, once more
  \item \emph{Say it:} say \emph{pautraḥ}
  \item \emph{Recall:} say the Sanskrit for coos, then the Sanskrit for chatters, then say \emph{pautraḥ} again
\end{itemize}

\begin{tcolorbox}[breakable,colback=teal!4,colframe=teal!35!black,title={Before you move on}]
What does \sk{पौत्रः} mean? (\textquotedblleft{}A grandson\textquotedblright{}.) Say it once more.
\end{tcolorbox}
\section[\texorpdfstring{\sk{पौत्री} (pautrī)}{पौत्री (pautrī)}]{\sk{पौत्री} (pautrī) — a granddaughter}

\label{lesson:SA-C146-pautri}



\begin{quote}
Before the new one: say the Sanskrit for chatters, then the Sanskrit for a grandson.
\end{quote}

\subsection*{You'll want to know: \sk{पौत्री}}
\textbf{\sk{पौत्री}} — \emph{pautrī} — \textquotedblleft{}a granddaughter\textquotedblright{}.

\begin{grammarlens}[title={What you've built}]
One more word for this chapter.
\end{grammarlens}

\subsection*{Your turn}
\begin{itemize}
\raggedright
  \item \emph{Say it:} \emph{pautrī}
  \item \emph{Say it:} \emph{pautrī}, once more
  \item \emph{Say it:} say \emph{pautrī}
  \item \emph{Recall:} say the Sanskrit for chatters, then the Sanskrit for a grandson, then say \emph{pautrī} again
\end{itemize}

\begin{tcolorbox}[breakable,colback=teal!4,colframe=teal!35!black,title={Before you move on}]
What does \sk{पौत्री} mean? (\textquotedblleft{}A granddaughter\textquotedblright{}.) Say it once more.
\end{tcolorbox}
\section[\texorpdfstring{\sk{पितामही} (pitāmahī)}{पितामही (pitāmahī)}]{\sk{पितामही} (pitāmahī) — a paternal grandmother}

\label{lesson:SA-C146-pitamahi}



\begin{quote}
Before the new one: say the Sanskrit for a grandson, then the Sanskrit for a granddaughter.
\end{quote}

\subsection*{You'll want to know: \sk{पितामही}}
\textbf{\sk{पितामही}} — \emph{pitāmahī} — \textquotedblleft{}a paternal grandmother\textquotedblright{}.

\begin{grammarlens}[title={What you've built}]
One more word for this chapter.
\end{grammarlens}

\subsection*{Your turn}
\begin{itemize}
\raggedright
  \item \emph{Say it:} \emph{pitāmahī}
  \item \emph{Say it:} \emph{pitāmahī}, once more
  \item \emph{Say it:} say \emph{pitāmahī}
  \item \emph{Recall:} say the Sanskrit for a grandson, then the Sanskrit for a granddaughter, then say \emph{pitāmahī} again
\end{itemize}

\begin{tcolorbox}[breakable,colback=teal!4,colframe=teal!35!black,title={Before you move on}]
What does \sk{पितामही} mean? (\textquotedblleft{}A paternal grandmother\textquotedblright{}.) Say it once more.
\end{tcolorbox}
\section[\texorpdfstring{\sk{मातामहः} (mātāmahaḥ)}{मातामहः (mātāmahaḥ)}]{\sk{मातामहः} (mātāmahaḥ) — a maternal grandfather}

\label{lesson:SA-C146-matamahah}



\begin{quote}
Before the new one: say the Sanskrit for a granddaughter, then the Sanskrit for a paternal grandmother.
\end{quote}

\subsection*{You'll want to know: \sk{मातामहः}}
\textbf{\sk{मातामहः}} — \emph{mātāmahaḥ} — \textquotedblleft{}a maternal grandfather\textquotedblright{}.

\begin{grammarlens}[title={What you've built}]
One more word for this chapter.
\end{grammarlens}

\subsection*{Your turn}
\begin{itemize}
\raggedright
  \item \emph{Say it:} \emph{mātāmahaḥ}
  \item \emph{Say it:} \emph{mātāmahaḥ}, once more
  \item \emph{Say it:} say \emph{mātāmahaḥ}
  \item \emph{Recall:} say the Sanskrit for a granddaughter, then the Sanskrit for a paternal grandmother, then say \emph{mātāmahaḥ} again
\end{itemize}

\begin{tcolorbox}[breakable,colback=teal!4,colframe=teal!35!black,title={Before you move on}]
What does \sk{मातामहः} mean? (\textquotedblleft{}A maternal grandfather\textquotedblright{}.) Say it once more.
\end{tcolorbox}
\section[\texorpdfstring{\sk{पितृव्यः} (pitṛvyaḥ)}{पितृव्यः (pitṛvyaḥ)}]{\sk{पितृव्यः} (pitṛvyaḥ) — a father's brother}

\label{lesson:SA-C146-pitrvyah}



\begin{quote}
Before the new one: say the Sanskrit for a paternal grandmother, then the Sanskrit for a maternal grandfather.
\end{quote}

\subsection*{You'll want to know: \sk{पितृव्यः}}
\textbf{\sk{पितृव्यः}} — \emph{pitṛvyaḥ} — \textquotedblleft{}a father's brother\textquotedblright{}.

\begin{grammarlens}[title={What you've built}]
One more word for this chapter.
\end{grammarlens}

\subsection*{Your turn}
\begin{itemize}
\raggedright
  \item \emph{Say it:} \emph{pitṛvyaḥ}
  \item \emph{Say it:} \emph{pitṛvyaḥ}, once more
  \item \emph{Say it:} say \emph{pitṛvyaḥ}
  \item \emph{Recall:} say the Sanskrit for a paternal grandmother, then the Sanskrit for a maternal grandfather, then say \emph{pitṛvyaḥ} again
\end{itemize}

\begin{tcolorbox}[breakable,colback=teal!4,colframe=teal!35!black,title={Before you move on}]
What does \sk{पितृव्यः} mean? (\textquotedblleft{}A father's brother\textquotedblright{}.) Say it once more.
\end{tcolorbox}
